\chapter{Luis A. Caffarelli (2023)}

\section{Biographical Sketch}



Luis \'Angel Caffarelli was born on 8 December 1948 in Buenos Aires, Argentina. He studied at the University of Buenos Aires, earning his undergraduate degree in 1969, and at the University of C\'ordoba, where he completed his PhD in 1972 under the supervision of Calixto Calder\'on. His early work on free boundary problems immediately established him as a leading analyst.

Caffarelli held positions at the University of Minnesota (1973--1980), the University of Chicago (1980--1986), the Institute for Advanced Study in Princeton (1986--1996), and the University of Texas at Austin (1996--present), where he holds the Sid W. Richardson Foundation Regents Chair in Mathematics.

He is one of the most honored analysts of his generation, receiving the Abel Prize in 2023, the Wolf Prize in Mathematics in 2012, the Shaw Prize in Mathematical Sciences in 2018, the Leroy P. Steele Prize for Lifetime Achievement in 2019, and the Rolf Schock Prize in 2005.

\section{High-Level View for a General Audience}

\subsection{What Did Caffarelli Do?}

Imagine a block of ice sitting in a warm room. The ice melts, and the boundary between ice and water moves over time. The shape of this moving boundary---the interface between solid and liquid---is not fixed in advance; it depends on the flow of heat in a complicated way.

This is a \emph{free boundary problem}: you need to find not just the temperature at each point, but also the location of the boundary itself. The problem is "free" because the boundary is part of the solution, not part of the problem statement.

Luis Caffarelli developed the mathematical theory for understanding the regularity of free boundaries---showing that under very general conditions, these boundaries are smooth surfaces (except possibly for a small set of singularities).

\subsection{Why Does It Matter?}

Caffarelli's work on the regularity of solutions to partial differential equations (PDEs) is essential for understanding:
\begin{itemize}
\item \textbf{Phase Transitions}: How ice melts, how water boils, and how metals solidify.
\item \textbf{Fluid Dynamics}: The behavior of the Navier--Stokes equations that describe fluid flow (Caffarelli proved the best-known regularity result for these equations).
\item \textbf{Optimal Transport}: The problem of moving mass from one distribution to another with minimum cost, which has applications from economics to machine learning.
\item \textbf{Materials Science}: The behavior of materials near cracks and defects.
\item \textbf{Image Processing}: Some of the most powerful techniques for denoising and enhancing images are based on PDEs that Caffarelli studied.
\end{itemize}

\subsection{The Big Picture}

Caffarelli's overarching contribution is to the \emph{regularity theory} of PDEs. When a physicist or engineer solves a differential equation on a computer, the solution is typically a smooth function. But mathematically, many PDEs can have solutions that are rough, singular, or even discontinuous. Caffarelli showed that for a very wide class of equations---including the most important ones from physics and geometry---the solutions are necessarily smooth, except for a small set of singularities that can be precisely characterized.


\begin{center}
\begin{tikzpicture}[scale=0.8]
  \draw[thick, abelgray] (-3,-2) rectangle (3,2);
  \draw[domain=-3:3, smooth, variable=\x, abelred, very thick] plot (\x, {0.5*sin(\x r)});
  \node[abelblue] at (-1.5, 1) {Liquid Phase};
  \node[abelgray] at (1.5, -1) {Solid Phase (Ice)};
  \node[abelred, right] at (1, 0.8) {Free Boundary $\Gamma$};
\end{tikzpicture}
\end{center}

\section{The Origin of the Problem}

\subsection{The Regularity Problem for Elliptic PDEs}

An elliptic partial differential equation is one that satisfies a certain positivity condition, analogous to the Laplace equation:
\[
\Delta u = \frac{\partial^2 u}{\partial x_1^2} + \cdots + \frac{\partial^2 u}{\partial x_n^2} = 0.
\]

The fundamental problem of regularity theory is: given a weak solution of an elliptic PDE, how smooth is it? For linear equations, the classical theory (Schauder estimates, Calder\'on--Zygmund theory) provides strong answers: solutions are $C^\infty$ if the coefficients are smooth.

\begin{knowledgebridge}{What Makes a PDE Elliptic?}
Elliptic equations are the ``steady-state'' cousins of wave and heat equations. Just as a stretched membrane finds a smooth equilibrium shape, elliptic PDEs tend to produce smooth solutions --- the key question is \emph{how} smooth, and under what conditions.
\end{knowledgebridge}

For nonlinear equations, the situation is far more delicate. The equation may have solutions that are not smooth, and understanding when and where smoothness fails is one of the central problems of analysis.

\subsection{Free Boundary Problems}

A free boundary problem is a PDE where part of the boundary is unknown and must be determined as part of the solution. The prototype is the Stefan problem (ice melting):
\[
u_t = \Delta u \quad \text{in the water region},
\]
\[
u = 0 \quad \text{on the free boundary},
\]
\[
V = -|\nabla u| \quad \text{(Stefan condition)}.
\]

Here $u$ is the temperature, and the free boundary separates the ice (where $u < 0$) from the water (where $u > 0$). The condition connecting the velocity $V$ of the interface to the temperature gradient makes this a highly nonlinear problem.

\subsection{The Obstacle Problem}

Another fundamental free boundary problem is the obstacle problem: find the equilibrium position of an elastic membrane pressed against an obstacle. This is equivalent to finding a function $u \geq 0$ such that:
\[
\Delta u \leq 0,\quad u \cdot \Delta u = 0.
\]

The free boundary is the boundary of the set where $u > 0$, separating the "free" part of the membrane from the part in contact with the obstacle.

\begin{primer}{The Obstacle Problem}
Picture a blanket draped over a basketball. Part of the blanket touches the ball (the ``obstacle''), and the rest hangs free. The boundary between the touching and free regions is the \emph{free boundary}. The obstacle problem asks: where does the blanket lose contact, and what shape does the free part take?
\end{primer}

\subsection{The Monge--Amp\`ere Equation}

The Monge--Amp\`ere equation is a fully nonlinear PDE:
\[
\det(D^2 u) = f \quad \text{in } \Omega,
\]
where $D^2 u$ is the Hessian matrix of $u$. This equation arises in:
\begin{itemize}
\item Optimal transport (Monge--Kantorovich problem).
\item Differential geometry (prescribed Gaussian curvature).
\item Fluid dynamics (Monge--Amp\`ere equation for potential flow).
\end{itemize}

\section{Deep Insight into the Problem}

\subsection{The Viscosity Solution Approach}

For fully nonlinear elliptic equations, classical solutions may not exist. Caffarelli, working with Evans, Krylov, and others, developed the theory of \emph{viscosity solutions}. A viscosity solution is defined using the comparison principle: $u$ is a viscosity subsolution if for any smooth test function $\phi$ that touches $u$ from above at a point $x_0$, we have:
\[
F(D^2\phi(x_0), D\phi(x_0), \phi(x_0), x_0) \leq 0.
\]

The viscosity framework allows for the existence and uniqueness of solutions under very mild conditions, and Caffarelli showed that these solutions are more regular than one might expect.

\subsection{The Evans--Krylov Theorem}

The Evans--Krylov theorem, which Caffarelli helped extend and generalize, states that viscosity solutions of uniformly elliptic fully nonlinear equations are $C^{2,\alpha}$ in the interior. This is the optimal regularity result for these equations and is essential for all further developments.

\subsection{Free Boundary Regularity}

Caffarelli's deep insight into free boundary problems was to classify the possible "blow-up" limits of the free boundary near a point. By rescaling the solution near a boundary point, one obtains a limiting solution that satisfies a simpler PDE. Caffarelli showed that for a wide class of free boundary problems, the possible blow-up limits are either:
\begin{itemize}
\item A flat free boundary (the half-space solution).
\item A corner-type solution (for singular points).
\end{itemize}

This classification leads to the regularity of the free boundary: the set of regular points (where the blow-up limit is flat) is an open set on which the free boundary is $C^\infty$, while the singular set has lower dimension.

\begin{bigidea}
Instead of examining the whole free boundary at once, Caffarelli's key idea was to zoom in on a single point --- like a microscope revealing the local structure. By studying these ``blow-up'' limits, he showed that only two patterns can appear: flat (smooth) or cusp-like (singular). This classification is the foundation of free boundary regularity.
\end{bigidea}

\section{Biggest Contributions and New Tools}

\subsection{The Caffarelli--Kohn--Nirenberg Theorem on the Navier--Stokes Equations}

This is one of the most important results in the mathematical theory of fluid dynamics. For the three-dimensional Navier--Stokes equations:
\[
u_t + u \cdot \nabla u - \Delta u + \nabla p = 0,\quad \nabla \cdot u = 0,
\]
the question of whether smooth solutions exist for all time is one of the Clay Millennium Problems.

Caffarelli, Kohn, and Nirenberg proved that for any weak solution satisfying a certain energy inequality, the singular set has zero one-dimensional Hausdorff measure. This means:
\begin{itemize}
\item The singular set is at most a set of isolated points in space-time.
\item The solution is smooth except possibly for a very small set of singularities.
\item This is the best regularity result currently known for the 3D Navier--Stokes equations.
\end{itemize}

\begin{analogy}{Navier--Stokes: The Ultimate Weather Forecast}
Weather prediction solves the Navier--Stokes equations for the atmosphere. The Caffarelli--Kohn--Nirenberg theorem guarantees that even when a forecast breaks down into turbulence, the breakdown is confined to isolated points --- not entire storms. It is the best guarantee of predictability we have.
\end{analogy}

\subsection{Free Boundary Regularity}

Caffarelli's theory of free boundary regularity includes:
\begin{itemize}
\item \textbf{The One-Phase Stefan Problem}: The free boundary is $C^\infty$ except on a set of measure zero.
\item \textbf{The Two-Phase Stefan Problem}: Regularity of the free boundary for the mushy region.
\item \textbf{The Obstacle Problem}: The free boundary is $C^{1,1}$ and the singular set has codimension at least 2.
\item \textbf{The Bernoulli Problem}: Regularity of the free boundary for the exterior problem.
\end{itemize}

\subsection{The Monge--Amp\`ere Equation}

Caffarelli's contributions to the Monge--Amp\`ere equation:
\begin{itemize}
\item \textbf{Interior $C^{1,\alpha}$ Regularity}: Strictly convex solutions of $\det(D^2 u) = f$ with $f$ bounded between positive constants are $C^{1,\alpha}$ in the interior.
\item \textbf{Interior $C^{2,\alpha}$ Regularity}: Under additional conditions, the solution is $C^{2,\alpha}$ in the interior.
\item \textbf{Global Regularity}: Regularity up to the boundary for strictly convex domains.
\item \textbf{The Caffarelli--Guti\'errez Theorem}: An $L^p$ theory for the linearized Monge--Amp\`ere equation.
\end{itemize}

\subsection{Fully Nonlinear Elliptic Equations}

Caffarelli's work on fully nonlinear equations includes:
\begin{itemize}
\item \textbf{The Evans--Krylov--Caffarelli Theorem}: $C^{2,\alpha}$ regularity for viscosity solutions of uniformly elliptic fully nonlinear equations.
\item \textbf{The Theory of Viscosity Solutions}: Existence, uniqueness, and stability of viscosity solutions.
\item \textbf{Aleksandrov--Bakelman--Pucci Estimates}: Sharp estimates for solutions of fully nonlinear equations.
\item \textbf{The Harnack Inequality}: An extension of the Harnack inequality to fully nonlinear equations.
\end{itemize}

\subsection{Nonlinear Diffusion and the Porous Medium Equation}

Caffarelli made fundamental contributions to the study of nonlinear diffusion equations:
\begin{itemize}
\item \textbf{The Porous Medium Equation}: $u_t = \Delta(u^m)$ for $m > 1$. Caffarelli proved regularity of solutions and the finite speed of propagation.
\item \textbf{The $p$-Laplacian}: Regularity of solutions to $\nabla \cdot (|\nabla u|^{p-2} \nabla u) = 0$.
\item \textbf{Fractional Diffusion}: Regularity for equations involving fractional Laplacians.
\end{itemize}

\section{Impact of the Discovery in Science and Technology}

\subsection{Impact on Mathematics}

Caffarelli's work has transformed the field of PDEs:

\begin{itemize}
\item \textbf{Regularity Theory}: The Evans--Krylov--Caffarelli theorem and Caffarelli's free boundary regularity are fundamental results taught in every graduate course on PDEs.

\item \textbf{Navier--Stokes Equations}: The Caffarelli--Kohn--Nirenberg theorem is the best regularity result for the 3D Navier--Stokes equations and remains an essential reference.

\item \textbf{Optimal Transport}: Caffarelli's regularity results for the Monge--Amp\`ere equation provide the theoretical foundation for the modern theory of optimal transport.

\item \textbf{Calculus of Variations}: Caffarelli's methods for free boundary problems are now standard tools in the calculus of variations.

\item \textbf{Geometric Analysis}: Caffarelli's work on fully nonlinear equations and the Monge--Amp\`ere equation is essential for geometric problems involving curvature.
\end{itemize}

\begin{whyitmatters}
Caffarelli's regularity theorems are the invisible scaffolding beneath modern applied mathematics. Every time an engineer simulates melting metal, a geophysicist models glacier flow, or a data scientist uses optimal transport to align datasets, Caffarelli's proofs assure them that the underlying equations behave as expected.
\end{whyitmatters}

\subsection{Impact on Science and Engineering}

\begin{itemize}
\item \textbf{Materials Science}: Free boundary problems model phase transitions in materials. Caffarelli's regularity results provide theoretical guarantees for numerical simulations of solidification and melting.

\item \textbf{Fluid Dynamics}: The Caffarelli--Kohn--Nirenberg theorem provides the best understanding of when and where singularities can form in fluid flows.

\item \textbf{Economics}: Optimal transport theory, built on Caffarelli's Monge--Amp\`ere regularity, is used in matching markets, resource allocation, and the theory of economic equilibrium.

\item \textbf{Machine Learning}: The theory of optimal transport is used in generative models (Wasserstein GANs), domain adaptation, and the analysis of fairness in algorithms.

\item \textbf{Image Processing}: Nonlinear diffusion equations of the type Caffarelli studied are used for image denoising, inpainting, and segmentation.

\item \textbf{Finance}: The Black--Scholes equation and its generalizations are PDEs of the type studied by Caffarelli.
\end{itemize}

\section{Current Developments}

\subsection{The Navier--Stokes Millennium Problem}

The Clay Millennium Problem for the Navier--Stokes equations remains open. The Caffarelli--Kohn--Nirenberg theorem provides the best current regularity result. Current research focuses on:
\begin{itemize}
\item Improving the dimension of the singular set.
\item Proving global regularity for special classes of initial data.
\item Connections with the theory of turbulence.
\end{itemize}

\subsection{Transportation Theory and Machine Learning}

Optimal transport is one of the most active areas of current research:
\begin{itemize}
\item Caffarelli's regularity results are essential for numerical methods for optimal transport.
\item The Wasserstein distance is used in machine learning for generative modeling.
\item The theory of entropic regularization of optimal transport builds on Caffarelli's work.
\end{itemize}

\subsection{Nonlocal Equations and the Fractional Laplacian}

Caffarelli's recent work on nonlocal equations (involving fractional Laplacians) has opened a new field:
\begin{itemize}
\item The theory of nonlocal minimal surfaces.
\item Regularity of solutions to integro-differential equations.
\item Connections to stochastic processes and L\'evy flights.
\end{itemize}

\subsection{Mean Field Games}

Caffarelli's work on fully nonlinear PDEs is essential for the theory of mean field games:
\begin{itemize}
\item The Hamilton--Jacobi--Bellman equations of mean field games are fully nonlinear.
\item Caffarelli's regularity results provide the foundation for the existence theory.
\end{itemize}

\section{Conclusion}

Luis Caffarelli's contributions to the regularity theory of nonlinear PDEs, free boundary problems, and the Monge--Amp\`ere equation are among the most important in modern analysis. His work has provided the theoretical foundations for understanding the smoothness of solutions to the equations that describe the physical world.

The Abel Prize citation recognized him "for his seminal contributions to regularity theory for nonlinear partial differential equations including free-boundary problems and the Monge--Amp\`ere equation."


\section{Behind the Breakthrough: The Human Story}
Caffarelli's physical intuition is legendary. Colleagues say he can 'see' the behavior of fluids and boundaries geometrically, allowing him to prove theorems that others find mathematically impenetrable.

\section{Pedagogical Metaphors and Lineage}
\subsection{Signature Metaphor}
``Predicting the exact shape of a melting block of ice as it turns into water.''

\subsection{Mathematical Lineage}
Advised by Calixto Calder\'on.

\section{Applied Science and Computational Algorithms}
\subsection{Key Takeaways for Applied Science}
Free boundary problems model physical phase transitions (melting ice), oil reservoir filtration (the Stefan problem), and pricing financial options (American options).

\subsection{Algorithms and Numerical Implementations}
Level Set methods (Osher-Sethian) and front-tracking finite element schemes numerically track moving boundaries and phase interfaces.

\section{The Research Frontier}
Regularity theory for non-local equations (fractional Laplacian operators) and applying viscosity solution techniques to non-linear turbulent flows.

\section{Guided Exercises and Challenges}
\begin{enumerate}[label=\textbf{\arabic*.}]
  \item Formulate the classical obstacle problem as a variational inequality and identify the free boundary.
  \item Describe the Caffarelli-Kohn-Nirenberg (CKN) theorem on the partial regularity of weak solutions to the Navier-Stokes equations.
  \item What is a viscosity solution to a second-order elliptic PDE? Why is this concept essential for the Monge-Amp\'ere equation?
\end{enumerate}

\section{Curriculum Map and Further Reading}
For students wishing to delve deeper into the mathematical machinery underlying these breakthroughs, the following learning path is recommended:
\begin{itemize}
  \item Luis Caffarelli and Sandro Salsa, \emph{A Geometric Approach to Free Boundary Problems}, AMS Graduate Studies.
  \item Luis Caffarelli and Xavier Cabr\'e, \emph{Fully Nonlinear Elliptic Equations}, AMS.
\end{itemize}

\section*{Bibliography}
\begin{enumerate}
\item L. A. Caffarelli, "The regularity of free boundaries in higher dimensions," Acta Mathematica 139 (1977), 155--184.
\item L. A. Caffarelli, R. Kohn, and L. Nirenberg, "Partial regularity of suitable weak solutions of the Navier--Stokes equations," Communications on Pure and Applied Mathematics 35 (1982), 771--831.
\item L. A. Caffarelli, "Interior $W^{2,p}$ estimates for solutions of the Monge--Amp\`ere equation," Annals of Mathematics 131 (1990), 135--150.
\item L. A. Caffarelli and X. Cabr\'e, "Fully Nonlinear Elliptic Equations," American Mathematical Society, 1995.
\item L. A. Caffarelli and C. E. Guti\'errez, "Properties of the solutions of the Monge--Amp\`ere equation," American Journal of Mathematics 119 (1997), 423--465.
\item L. A. Caffarelli and S. Salsa, "A Geometric Approach to Free Boundary Problems," American Mathematical Society, 2005.
\item L. A. Caffarelli and P. E. Souganidis, "A rate of convergence for monotone finite difference approximations to fully nonlinear, uniformly elliptic PDEs," Communications on Pure and Applied Mathematics 61 (2008), 1--17.
\item L. A. Caffarelli, "Free boundary problems," in "Proceedings of the International Congress of Mathematicians 1986."
\item G. De Philippis and A. Figalli, "The Monge--Amp\`ere equation and its link to optimal transportation," Bulletin of the American Mathematical Society 51 (2014), 527--580.
\item C. Villani, "Optimal Transport: Old and New," Springer, 2009.
\item A. Figalli, "Free Boundary Problems," in "Proceedings of the International Congress of Mathematicians 2022."
\item M. G. Crandall, H. Ishii, and P.-L. Lions, "User's guide to viscosity solutions of second order partial differential equations," Bulletin of the American Mathematical Society 27 (1992), 1--67.
\item J. L. V\'azquez, "The Porous Medium Equation," Oxford University Press, 2007.
\item L. Silvestre, "Regularity of the obstacle problem for a fractional power of the Laplace operator," Communications on Pure and Applied Mathematics 60 (2007), 67--112.
\item T. Tao, "Nonlinear dispersive equations: local and global analysis," American Mathematical Society, 2006.
\end{enumerate}
